%% ============================================================================
%% Fichier  : annexes/Ann_E_bibliographie.tex
%% Titre    : Bibliographie Commentée
%% Version  : 1.0 -- novembre 2026 (INF3226, d'après ACSI)
%% ============================================================================
%% Uniquement des entrées de bibliographie/ASI_references.bib ; seuls A0 et A1
%% (HAL) sont cités parmi les textes de l'auteur.
%% ============================================================================

\chapter{Bibliographie Commentée}
\label{ann:bibliographie}

Pour chaque partie, quelques lectures, avec ce qu'elles apportent. Les
références complètes figurent dans la bibliographie du livre. Les textes
de l'auteur sont des prépublications déposées sur HAL~; aucun travail non
publié n'est cité.

\section*{Pour commencer}

\begin{itemize}
    \item \textbf{Le programme de l'auditologie}~\cite{minka2026auditologie}~:
          le texte de référence du cours (objet, axiomes, piliers,
          thématiques, méthode). À lire après le
          Chapitre~\ref{chap:audit-auditologie}, puis à relire avec le
          Chapitre~\ref{chap:auditologie}.
    \item \textbf{L'instrument trichotomique}~\cite{minka2026nek}~: la
          réponse \NEK{}, ses conditions, les métriques et le SUI~;
          fondement des Chapitres~\ref{chap:preuve}, \ref{chap:conclure}
          et~\ref{chap:incertitude}.
\end{itemize}

\section*{Partie I · Fondements}

\begin{itemize}
    \item COSO, \emph{Internal Control -- Integrated Framework}~\cite{coso2013ic}~:
          les cinq composantes et dix-sept principes~; la référence du
          contrôle interne.
    \item ISO~31000~\cite{iso31000_2018} et ISO/IEC~27005~\cite{iso27005_2022}~:
          la gestion des risques, générale puis appliquée à la sécurité de
          l'information (édition 2022, qui remplace celle de 2018).
    \item Les méthodes d'appréciation~: EBIOS Risk
          Manager~\cite{anssi2018ebios}, MEHARI~\cite{clusif2010mehari},
          OCTAVE Allegro~\cite{caralli2007octave}, FAIR~\cite{freund2014fair}~;
          Hubbard et Seiersen~\cite{hubbard2016cyber} pour la mesure
          quantitative, Cox~\cite{cox2008matrices} pour les limites des
          matrices de risque.
    \item Le guide COSO sur le risque de fraude~\cite{coso2016fraude} et
          l'étude de Cressey~\cite{cressey1953}, à l'origine du triangle de
          la fraude~; le diamant de Wolfe et Hermanson~\cite{wolfe2004diamond}~;
          Benford~\cite{benford1938} et Nigrini~\cite{nigrini2012benford} pour
          l'analyse des premiers chiffres.
    \item IIA, le modèle des trois lignes~\cite{iia2020troislignes}.
\end{itemize}

\section*{Partie II · Méthode}

\begin{itemize}
    \item ISO~19011~\cite{iso19011_2018}~: les sept principes de l'audit et
          la conduite d'un programme d'audit.
    \item INTOSAI, GUID~5100~\cite{intosai2019guid5100}~: la vérification des
          SI dans le secteur public, en français.
    \item Miller~\cite{miller1956}~: l'article sur la mémoire de travail qui
          fonde la borne de sept.
    \item Shannon~\cite{shannon1948}, Kullback et Leibler~\cite{kullback1951}~:
          l'entropie et la divergence, aux sources des mesures du cours~; Cover et
          Thomas~\cite{cover2006elements} pour les démonstrations.
    \item Simon~\cite{simon1962complexity}~: la quasi-décomposabilité des
          systèmes complexes~; Minto~\cite{minto2009pyramid}~: le principe
          pyramidal, utile à la décomposition et à la rédaction.
\end{itemize}

\section*{Partie III · Cadre normatif, juridique et déontologique}

\begin{itemize}
    \item IIA, Normes mondiales d'audit interne~\cite{iia2024gias}, et le
          CRIPP de 2017~\cite{iia2017cripp}.
    \item ISACA, ITAF~\cite{isaca2020itaf}~; IAASB, ISA~315~\cite{iaasb2019isa315}
          et ISAE~3402~\cite{iaasb2011isae3402}.
    \item IDI, manuel de l'audit de conformité~\cite{idi2021conformite}~:
          la relation tripartite, les critères, les niveaux d'assurance.
    \item COBIT~2019~\cite{isaca2018cobitintro,isaca2018cobitobjectifs} et
          COBIT~4.1~\cite{itgi2007cobit41}~; ISO/IEC~27001 et 27002 (édition
          2022)~\cite{iso27001_2022,iso27002_2022}~; NIST SP~800-53 et
          800-53A~\cite{nist2020sp80053,nist2022sp80053a}~; guide du
          CHAI~\cite{chai2014guide}.
    \item Loi n°~2010/012~\cite{loi2010012}~; règlement COBAC
          R-2016/04~\cite{cobac2016r04}~; AUSCGIE~\cite{ohada2014auscgie}~;
          régime financier de l'État~\cite{minfi2023regimefinancier}.
    \item NIST CSF~2.0~\cite{nist2024csf}~; code d'éthique de
          l'ISACA~\cite{isaca2024ethique} et code de l'IESBA~\cite{iesba2018code}
          pour la déontologie.
\end{itemize}

\section*{Partie IV · Pratique}

\begin{itemize}
    \item Le guide du CHAI~\cite{chai2014guide}, fiches 3.1 à 3.8~: des
          points de contrôle par thème, à adapter à chaque mission.
    \item ISO/IEC~27007~\cite{iso27007_2020} pour l'audit d'un SMSI~;
          ISO~22301~\cite{iso22301_2019} pour la continuité.
    \item CIS Controls~\cite{cis2021controls}, CVSS~\cite{first2019cvss},
          NIST SP~800-115~\cite{nist2008sp800115}, OWASP Top~10 et
          ASVS~\cite{owasp2021top10,owasp2021asvs}~: la sécurité technique et
          les tests d'intrusion.
    \item NIST SP~800-145~\cite{nist2011sp800145}, ISO/IEC~27017 et
          27018~\cite{iso27017_2015,iso27018_2019}, CSA
          CCM~\cite{csa2021ccm}~: le nuage~; ISO~37301~\cite{iso37301_2021}~:
          les systèmes de management de la conformité.
    \item La décision de la SEC sur Knight Capital~\cite{sec2013knight}~: un
          cas réel de changement mal maîtrisé, utile pour les
          Chapitres~\ref{chap:audit-projets} et~\ref{chap:audit-continuite}.
\end{itemize}

\section*{Partie V · Défis avancés}

\begin{itemize}
    \item NIST AI RMF~\cite{nist2023airmf}, ISO/IEC~42001~\cite{iso42001_2023}
          et la Recommandation de l'UNESCO~\cite{unesco2021ethique}~: auditer
          un système d'IA~; le règlement européen sur
          l'IA~\cite{ue2024aiact}, comme référentiel volontaire.
    \item Wainer et al.~\cite{wainer2000cat}~: les tests adaptatifs, aux
          sources des grilles adaptatives.
    \item ISO/IEC~27037~\cite{iso27037_2012}~: la preuve numérique et la
          chaîne de conservation, prolongées par les UE SEC4031 et SEC4052.
\end{itemize}

\section*{Partie VI · Recherche}

\begin{itemize}
    \item Hevner et al.~\cite{hevner2004design}~: la science du design en
          systèmes d'information.
    \item Lakatos~\cite{lakatos1970}~: noyau dur et ceinture protectrice.
    \item Simon~\cite{simon1957} et Kahneman~\cite{kahneman2011}~:
          rationalité limitée et biais~; von~Neumann et
          Morgenstern~\cite{vonneumann1944}, Nash~\cite{nash1951},
          Harsanyi~\cite{harsanyi1967}~: les jeux~; Norris~\cite{norris1998}~:
          les chaînes de Markov~; Bhaskar~\cite{bhaskar1975},
          Simondon~\cite{simondon1958}, Latour~\cite{latour2005}~: l'ontologie.
    \item Fellingham et Newman~\cite{fellingham1985strategic}~: l'audit
          comme jeu~; Spence~\cite{spence1973signaling}~: le signal
          coûteux~; le GUM~\cite{jcgm2008gum}~: l'expression de l'incertitude
          de mesure.
\end{itemize}
